\subsection{A shared Fermat-twelve quotient and both actual cover matrices}

\begin{statement}{D.37}{Proved}{A generating mixed Fermat shadow with exact elliptic-cover evaluations}
Retain $B,F,Q,N,x,y,t$ from D.25, and put
\[
 \begin{gathered}
 F_{12}=\ker\bigl(F\longrightarrow(\mathbb Z/12)^2\bigr),
 \qquad L=N\cap F_{12},\\
 H=B/L,\qquad J_m=B/[L,L]L^m .
 \end{gathered}                                           \tag{D.37.1}
\]
The character sends $x,y$ to the indicated ordered basis. Then
\[
 \begin{gathered}
 |H|=6912,\qquad \operatorname{rank}L=1153,\\
 \ker(H\to Q)\simeq C_2^2\times C_3,\qquad
 |J_m|=6912m^{1153}.
 \end{gathered}                                           \tag{D.37.2}
\]
There is an explicitly certified automorphism of $J_{12}$ with
multiplier one and the extra CM pair
\[
 A_{12}=I,\qquad
 \beta_{12}=\begin{pmatrix}9&8\\4&9\end{pmatrix},\qquad
 V_{12}=I.                                               \tag{D.37.3}
\]
Before the common Ihara gauge, its two conjugators project to
exactly the retained $G_6$ conjugators of D.33: the actual
pro-two component, including its right-$s$-centralizer correction,
and the separately specified pro-three component. The constructed
$J_{12}$ action also retains its recorded $J_6$ action. Agreement
with D.33's larger degree-eight quotient is not asserted.

The cusp has order $288$ in $J_{12}$. Its pullback to $B_3$
therefore has $N_{\rm ord}=144$. The recorded gauge makes $f$ a
commutator word and gives both $m=0$ hexagons and generation,
so (D.37.3) occurs in an actual nonarithmetical finite GT0 shadow.

Restriction to the three marked cyclic covers of D.36 and
filling their peripheral classes give actual maps from this
finite source. In the bases recorded there, the normalized
shadow has
\[
 \rho_{C_3,4}=-I,\qquad \rho_{C_4,3}=I.                  \tag{D.37.4}
\]
Both matrices are reductions of its one compact action on
$C_{12}:v^{12}=z(z-1)$, through the two explicit integral maps
of D.36. Its cubic Teichm\"uller factor is one. Consequently
its normalized cubic sign is $-1$ and its quartic projective
$i$-exponent is zero; these two actual cover coordinates violate
the arithmetic graph of D.36. Its retained fixed-fibre
coordinates remain $T_4=I$ and $W_3=[i]$.

The maximal common quotients with the full relative pro-two
and pro-three sources over $Q$ are also explicit. Let $H_4,H_3$
be the congruence quotient joined with the affine Fermat
characters modulo $4,3$, and let $L_j=\ker(B\to H_j)$.
Then $|H_4|=2304$, $|H_3|=1728$, and
\[
 \begin{aligned}
 J_{12}^{\langle2\rangle}&=B/[L_4,L_4]L_4^4,
       &\operatorname{rank}L_4&=385,\\
 J_{12}^{\langle3\rangle}&=B/[L_3,L_3]L_3^3,
       &\operatorname{rank}L_3&=289.
 \end{aligned}                                           \tag{D.37.5}
\]
Here the superscript means the universal quotient whose
kernel over $Q$ is a $p$-group. Their maximal common quotient
is $Q$, and the canonical map has the exact kernel
\[
 \begin{gathered}
 \ker\left(J_{12}\longrightarrow
 J_{12}^{\langle2\rangle}\mathbin{\times_Q}
 J_{12}^{\langle3\rangle}\right)\\
 \simeq (\mathbb Z/4)^{768}\oplus(\mathbb Z/3)^{864}.
 \end{gathered}
                                                               \tag{D.37.6}
\]
These maps compute the mixed part of this finite test. They
do not decide its lift to the full profinite group or the
attainability of the entire CM pair.

The manuscript passage addressed is Part D(e), section
51.16, folio 734 verso, scan 174, together with
D-after-LM2(1), section 51.15, folios 725--727,
scans 156--160. The comparison keeps the two exact cover
maps and the fixed-fibre maps of D.36.
\end{statement}
\begin{proof}
We first give the quotient maps, including the affine marking.
The following matrices and translations define a homomorphism
from $B$ to $\mathbb Z^2\rtimes S_3$:
\[
 s\longmapsto\left(
 \begin{pmatrix}0&1\\1&0\end{pmatrix},0\right),\qquad
 r\longmapsto\left(
 \begin{pmatrix}0&-1\\1&-1\end{pmatrix},e_1\right).
                                                               \tag{D.37.7}
\]
Both defining torsion relations hold. The linear part is the
specified $S_3$ quotient. Direct multiplication gives
$t^2$ the translation $-e_1$, so $x=t^{-2}$ and $y=sxs$
give $e_1,e_2$. The restriction to $F$ is therefore its
ordered abelianization, with kernel $F'$. Its image contains
every integral translation and both linear generators;
thus (D.37.7) identifies $B/F'$ with the displayed affine
group by actual inverse generators. Reducing its translations
modulo 12 and joining it with $Q$ gives precisely $B/L$.

An element of the joint image is recorded as $(q,a,b)$,
where $q\in Q$ and $(a,b)\in(\mathbb Z/12)^2$; multiplication
uses the linear matrix determined by $q\bmod2$.
Exact enumeration from the two displayed generators gives
6912 elements, with
\[
 \ker(H\to Q)=
 \{(a,b):a,b\equiv0\pmod2,\ a+b\equiv0\pmod3\}.
                                                               \tag{D.37.8}
\]
This is $C_2^2\times C_3$. Reduction of the same recorded
states modulo $4$ and modulo $3$ gives orders $2304$ and
$1728$; their fibre product over $Q$ has order 6912, so
the reduction map identifies that fibre product with $H$.
The cusp $t$ has order 24. The quotient maps to the full
Fermat-twelve affine group and to $Q$ have kernels of
orders 8 and 12, respectively.
The kernel in (D.37.8) is abelian of exponent six, so
$[N,N]N^6\subset L$. Conversely,
$[L,L]L^{12}\subset[N,N]N^6$. These inclusions give the
actual comparison chain
\[
 J_{12}\longrightarrow G_6\longrightarrow H\longrightarrow Q.
\]

The kernel $L$ is torsion-free since it lies in $F$.
The graph $L\backslash\mathcal T$ has 6912 edges and
$6912/2+6912/3=5760$ vertices. Its spanning tree gives
1153 fundamental cycles and free generators. In particular,
all ranks in (D.37.2) and (D.37.5) follow from these actual
covering graphs. The full Fermat subgroup $F_{12}$ has
rank 145, and $[F_{12}:L]=8$, which also checks the same
rank by the free-group covering formula.

Write $e_h$ for the oriented edge indexed by $h\in H$.
The path cocycle is
\[
 P(s)=0,\qquad P(r)=e_1-e_r,\qquad
 P(ab)=P(a)+aP(b).                                      \tag{D.37.9}
\]
Its restriction to $L$ is the integral cycle group. The
recorded cycle matrix $C$ has an identity submatrix on
its 1153 non-tree edges. Thus reduction modulo any $m$
faithfully represents $J_m$ as the image in
$(\mathbb Z/m)[H]\rtimes H$: its kernel on $L$ is exactly
$[L,L]L^m$, including for the nonfield coefficient rings
$\mathbb Z/4$ and $\mathbb Z/12$.

There are six possible $Q$ conjugators for $r$ and eight
for $s$ in D.25's normalized matrix problem. Each has
twelve lifts to $H$. Of the resulting $48\cdot144$
ordered choices, 576 satisfy
\[
 \beta s\beta^{-1}A rA^{-1}=t                           \tag{D.37.10}
\]
already in $H$, and each induces an automorphism there.
The code verifies this last assertion on all 6912 states
and both generating transitions. It records the induced
crossed homomorphism on the six linear $S_3$ elements;
the resulting affine map is a permutation.

Put $s'=\beta s\beta^{-1}$ and $r'=ArA^{-1}$ in $H$.
For corrections $b,a$ in the cycle lattice, the complete
cusp defect is the affine linear expression
\[
 \begin{split}
 \mathcal E(a,b)&=E+(1-s')b+s'(1-r')a,\\
 E={}&P_\beta-s'P_\beta\\
     &+s'\bigl(P_A-r'P_A+A P(r)\bigr)-P(t).
 \end{split}                                             \tag{D.37.11}
\]
Every term is evaluated in the exact graph module. Since
$s',r'$ generate $H$, the correction image has the same
coinvariants as the original $s,r$ image. Both its mod-two
and mod-three ranks are 1152, with one-dimensional
cokernel. Exhausting all 576 base choices gives
\[
\begin{array}{c|rrr}
 \text{linear cusp condition}&
 \text{base choices}&\text{matrix pairs}&\text{extra pairs}\\\hline
 \bmod2&288&24&12\\
 \bmod3&192&48&24\\
 \text{both, with the same base choice}&96&24&12.
\end{array}                                               \tag{D.37.12}
\]
These are linear solvability counts. Generation below is
verified for the specified final witness. The cocycle of
$t^{72}$ has zero image in each of the two cokernels, so
allowing the compatible lower-level multiplier changes
$\lambda=1+72j$ does not alter this table. The witness uses
$\lambda=1$ throughout.

We retain more than the two matrix coordinates. Under
$H\to Q$, the graph-chain map sends $e_h$ to $e_{q(h)}$.
On cycle coordinates this gives an explicit integral
$97\times1153$ matrix $S$. Its mod-two rank is 95. The
target mod-two conjugators are the D.27 homology masks
used throughout D.33, with the recorded right-centralizer
term on $\beta$. Every later D.33 correction lies in
$\Phi(N^{(2)})$, so this is its exact $G_2$ component.
For the separate odd component, the target is $A=1$ and
the D.25 path to $\beta$ plus twice its cycle number 39;
this is D.33's specified $G_3$ centralizer representative.
The proof uses that specified odd component, rather than
reducing D.33's discrete pro-two word modulo three.

Appending the two exact conjugator projections to
(D.37.11) gives ranks 1247 and 1248 at two and three,
and affine dimensions 1059 and 1058. The selected base
conjugators have affine coordinates
\[
 A=(I,(6,0)),\qquad
 \beta=\left(\begin{pmatrix}3&4\\8&3\end{pmatrix},(6,0)\right),
                                                               \tag{D.37.13}
\]
where the matrix is the canonical global-sign representative
of (D.37.3). The generated mod-two witness is the first
recorded choice; at three the successive kernel ranks are
$1099,1045,1135,1063,1081,1135,1153$. The seventh choice
is therefore the recorded generating one.

Here generation is an exact group assertion. If $\alpha_H$
is the verified base automorphism, the graph map sends
\[
 e_h\longmapsto \alpha_H(h)
       \bigl(P_A+a-P_\beta-b+e_A\bigr).                  \tag{D.37.14}
\]
Its restrictions to all 1153 fundamental cycles have rank
1153 at both primes. The induced endomorphism of $J_6$
is already an automorphism of $H$ and is invertible on its
kernel, so it is an automorphism. Every cycle is also
converted to a literal $B$ word and checked integrally.
Independent substitution of the two generator images into
all those words reproduces the entire kernel matrices.

The exponent-four step retains this very $J_6$ action.
Let $F_0$ be the integral correction matrix in (D.37.11),
and let $z$ be the chosen mod-two correction. The upper
equation is
\[
 F_0 w=-\frac{E+F_0z}{2}\pmod2,\qquad z_4=z+2w.           \tag{D.37.15}
\]
The numerator is even in every cycle coordinate. The
recorded witness has zero upper obstruction on its first
test. This equation has rank 1152 and affine dimension
1154. The resulting $z_4$ is combined with the same
mod-three solution by $z_{12}=9z_4+4z_3\pmod{12}$.
Invertibility modulo two makes the kernel matrix invertible
modulo four; together with the verified mod-three rank it
is invertible over $\mathbb Z/12$.

The final cusp word has 673392 letters in $s,r,r^{-1}$.
It evaluates to zero in all 6912 edge coordinates modulo
12, and all 1153 basis-image checks are repeated on the
final pair. The class of $t^{24}$ has a nonzero reduction
at both primes; its mod-twelve order is 12. This proves
the asserted cusp order 288.

For clarity, the finite GT normalization is also checked
on the same pair. Its integral $F$ abelianization is
$\beta^{\rm ab}=882x+180y$. Set $k=-702$ and
$\beta'=x^k\beta$. Then
\[
 f'=s\beta's(\beta')^{-1}\in F',\qquad
 \alpha'=\operatorname{Int}(x^k)\alpha.                 \tag{D.37.16}
\]
The exact abelianization of $f'$ is $(0,0)$. The gauge fixes
$u=t^{-1}$ and gives
$\alpha'(s)=sf'$, $\alpha'(r)=(f')^{-1}r$ and
$\alpha'(v)=(f')^{-1}vf'$, as in D.25. Both corresponding
$m=0$ hexagons are multiplied and checked in the faithful
$J_{12}$ model. The pullback of $[L,L]L^{12}$ to $B_3$
is normal, finite index, and contained in $PB_3$.
The braid center is killed, and $x=t^{-2}$ has order144.
Thus the commutator, unit, two-hexagon and surjectivity
conditions of \cite[Definitions 2.1 and 2.3]{dBDHP2026}
all hold for this actual finite shadow. The same left
gauge on $A,\beta,V$ leaves $T=A^{-1}V$ and
$W=\beta^{-1}V$ unchanged. Its projection to $G_6$ retains
D.25's conjugator-rigidity test, proving that (D.37.3)
is an invariant of this shadow and has no Galois realization.

We next prove (D.37.4) through the actual cover maps.
For $n=3,4,12$ put
$K_n=\ker(F\to\mathbb Z/n)$, with both loops mapping to one.
There are inclusions $L\subset F_{12}\subset K_{12}\subset K_n$.
The geometric Fermat marking also has an explicit map
from the manuscript's plus model. Over $\mathbb Q(\zeta_{24})$,
on $x_0^{12}+x_1^{12}+x_\infty^{12}=0$ it is
\[
 z=-\frac{x_0^{12}}{x_\infty^{12}},\qquad
 v=\zeta_{24}\frac{x_0x_1}{x_\infty^2},
 \qquad v^{12}=z(z-1).
\]
With $x_\infty=1$, choose the positive-interval lifts
$x_0=\zeta_{24}z^{1/12}$ and
$x_1=\zeta_{24}^{-1}(1-z)^{1/12}$; this gives exactly
D.36's cyclic-twelve tangent $v\sim\zeta_{24}z^{1/12}$.
The maps to $C_3,C_4$ are $v\mapsto v^4,v^3$.
Thus the even-degree twist and the field of this map
are retained. The finite multiplier here is one modulo288,
so every Galois element inducing this shadow would fix
that field.
For $m\mid12$,
\[
 [L,L]L^{12}\subset[K_n,K_n]K_n^m.                       \tag{D.37.17}
\]
Consequently restriction of the finite automorphism and
abelianization give the actual action on
$H_1(K_n,\mathbb Z/m)$. The ordered peripheral classes
are preserved downstairs; their lifted conjugacy classes
are permuted with the same powers. Filling their span
therefore gives the compact action. The code verifies
this span preservation directly as well.

There is no ambiguous free-basis conversion in this
calculation. Collapsing the tree of the six-edge $S_3$
graph gives chord generators $b_1,b_2$ with
$x=b_1$, $y=b_2^{-1}$, checked in both directions.
The actual normalized $f'$ becomes a 99308-letter word
in $x^{\pm1},y^{\pm1}$. For the Schreier basis
\[
 a=x^n,\qquad b_j=x^j yx^{-((j+1)\bmod n)}
                 \quad(0\le j<n),
\]
let $D_n$ denote conjugation by $x$ on punctured homology,
and let $v_n$ be the exact homology class of $f'$.
Since $\alpha'(x)=x$ and
$\alpha'(y)=(f')^{-1}yf'$, its homology matrix has columns
\[
 e_a,\qquad e_{b_j}+D_n^j(D_n-I)v_n.                      \tag{D.37.18}
\]
The program computes this matrix and independently
substitutes $x$ and $(f')^{-1}yf'$ into every Schreier word.
The integral filling matrices $P_n,J_n$ from D.36 give
$P_n H_nJ_n$. These are exactly $-I$ modulo four at
$n=3$ and $I$ modulo three at $n=4$.
The two equations
\[
 M_n H_{12}^{\rm c}=H_n^{\rm c}M_n
       \pmod{m_n},\qquad (n,m_n)=(3,4),(4,3),             \tag{D.37.19}
\]
are checked using the same compact cyclic-twelve action
and the actual matrices $M_n$ of D.36. Thus these are
evaluations of one source action under the specified maps.

The order-three lift of the cubic mod-two matrix is $I$,
so its normalized principal-unit sign is $-1$. The
quartic matrix is $I$, giving projective exponent zero.
For an actual Galois element with this finite multiplier,
D.36's tangent and cubic-twist comparisons impose its
six-element arithmetic graph, where these two signs
agree. Their computed disagreement gives a second direct
certificate of nonarithmeticality. It supplies no
unproved cover-to-fixed-fibre equality for GT0: both
maps have instead been evaluated on the same finite
shadow, giving the four coordinates stated above.

Finally, the universal quotient statements are exact
group calculations. The two subgroup sequences are
\[
 1\to L\to L_4\to C_3\to1,\qquad
 1\to L\to L_3\to C_2^2\to1.
\]
Their homology five-term sequences give
\[
 (L^{\rm ab}/4)_{C_3}\simeq L_4^{\rm ab}/4,\qquad
 (L^{\rm ab}/3)_{C_2^2}\simeq L_3^{\rm ab}/3,             \tag{D.37.20}
\]
because the complementary group order is a unit in the
coefficient ring. Equivalently, averaging its finite
action gives the same explicit projection on each
module. In a relative pro-two quotient of $J_{12}$,
the three-primary module and the normal closure of that
$C_3$ must be killed. The remaining module is exactly
the first coinvariant module in (D.37.20); its extension
over $H_4$ is $B/[L_4,L_4]L_4^4$. This extension has a
two-group kernel over $Q$, so it is itself a quotient
of the full relative pro-two source. This proves both
factorizations and universality. The same argument with
$C_2^2$ proves the pro-three assertion.

Any common quotient of those two relative quotients
has a kernel over $Q$ that is both a two-group and a
three-group, hence is trivial. They therefore have
maximal common quotient $Q$. Goursat's argument applied
to their common surjection from $J_{12}$ gives the full
fibre product in (D.37.6). On the coefficient kernel,
averaging over $C_3$ and over $C_2^2$ splits off modules
of ranks 385 and 289. The complementary ranks are
$1153-385=768$ and $1153-289=864$. Since a direct summand
of a finite free module over $\mathbb Z/4$ is free,
these complementary modules are exactly those in
(D.37.6). This proves the stated kernel, not only its
order. The existence of a lift of the recorded action
through the remaining full source is still the exact
all-finite-quotient question of D.22.
\end{proof}
